\documentclass[aspectratio=169]{beamer}
\usetheme{Madrid}
\usecolortheme{default}
\usepackage{amsmath}
\usepackage{graphicx}
\usepackage{tikz-cd}
\usepackage{multicol}

\setlength{\parindent}{0pt}
\setlength{\parskip}{1\baselineskip}

\title[Lecture 09]{An Introduction to E-Commerce\\\small with a Focus on Machine Learning, Deep Learning, and Artificial Intelligence}
\subtitle{Lecture 09: Pricing and Promotions with ML + Optimization}
\author{Dr. Haitham A. El-Ghareeb}
\institute{Faculty of Computers and Information Sciences \\ Mansoura University \\ Egypt}
\date{\today}

\begin{document}

\begin{frame}
  \titlepage
\end{frame}

\begin{frame}{Session plan (90 minutes)}
  \begin{enumerate}
    \item Why pricing is a core AI problem in commerce
    \item Pricing objectives
    \item Price, promotion, and discount structure
    \item Demand forecasting
    \item Price elasticity
    \item Hands-on lab: forecast plus constrained optimization
  \end{enumerate}
\end{frame}

\begin{frame}{Learning objectives}
  By the end of this lecture, you should be able to:
  \begin{itemize}
    \item Explain why pricing in e-commerce is both a prediction problem and an optimization problem.
    \item Distinguish demand forecasting, price elasticity estimation, and promotion uplift modeling.
    \item Analyze how business constraints such as inventory, margins, contracts, and competition shape pricing decisions.
    \item Design a pricing workflow that combines forecasting, optimization, and experimentation rather than treating price as a static input.
  \end{itemize}
\end{frame}

\section{Why pricing is a core AI problem in commerce}

\begin{frame}{Why pricing is a core AI problem in commerce}
  \begin{itemize}
    \item demand estimation,
    \item elasticity understanding,
    \item business-rule enforcement,
    \item optimization under constraints,
    \item and experimentation to measure causal impact.
  \end{itemize}
\end{frame}

\section{Pricing objectives}

\begin{frame}{Pricing objectives}
  \begin{itemize}
    \item \textbf{Common objectives:} \textbf{Revenue maximization:} increase gross sales volume.; \textbf{Margin or contribution optimization:} emphasize profitability rather than top-line revenue.
    \item \textbf{Why objective clarity matters:} clearing inventory,; preserving premium brand position,
  \end{itemize}
\end{frame}

\section{Price, promotion, and discount structure}

\begin{frame}{Price, promotion, and discount structure}
  \begin{itemize}
    \item \textbf{Promotions as structured interventions:} percentage discount,; fixed-amount coupon,
  \end{itemize}
\end{frame}

\section{Demand forecasting}

\begin{frame}{Demand forecasting}
  \begin{itemize}
    \item \textbf{What is being forecast:} unit demand by SKU and day,; category demand,
    \item \textbf{Inputs to demand models:} historical sales,; current and historical prices,
    \item \textbf{Why forecasting is not enough:} Forecasting predicts what may happen.
  \end{itemize}
\end{frame}

\section{Price elasticity}

\begin{frame}{Price elasticity}
  \begin{itemize}
    \item \textbf{Intuition:} If a small increase in price causes a large reduction in demand, the product is relatively price-sensitive.
    \item \textbf{Why elasticity is difficult to estimate:} In live commerce systems, observed prices are not randomly assigned.
    \item \textbf{Practical considerations:} category,; season,
  \end{itemize}
\end{frame}

\section{Promotions and uplift}

\begin{frame}{Promotions and uplift}
  \begin{itemize}
    \item \textbf{Observed uplift versus causal uplift:} genuine causal impact from the promotion,; seasonal demand shifts,
    \item \textbf{Why uplift matters:} Promotions can be expensive.
  \end{itemize}
\end{frame}

\section{Optimization under business constraints}

\begin{frame}{Optimization under business constraints}
  \begin{itemize}
    \item \textbf{Typical constraints:} minimum margin thresholds,; inventory availability,
    \item \textbf{Why unconstrained optimization is unrealistic:} violates margin policy,; creates stockouts,
  \end{itemize}
\end{frame}

\section{Dynamic pricing}

\begin{frame}{Dynamic pricing}
  \begin{itemize}
    \item \textbf{When dynamic pricing can help:} demand fluctuates quickly,; inventory is perishable or time-sensitive,
    \item \textbf{When dynamic pricing can hurt:} customer trust may decline if prices feel unstable or unfair,; operations may struggle with frequent updates,
  \end{itemize}
\end{frame}

\section{Competition-aware pricing}

\begin{frame}{Competition-aware pricing}
  \begin{itemize}
    \item \textbf{Competitive signals:} competitor list price,; promotion intensity,
    \item \textbf{Strategic caution:} A competitor may use a loss leader or temporary campaign that should not be copied directly.
  \end{itemize}
\end{frame}

\section{B2C versus B2B pricing}

\begin{frame}{B2C versus B2B pricing}
  \begin{itemize}
    \item \textbf{B2C:} high-volume SKU-level decisions,; campaign responsiveness,
    \item \textbf{B2B:} negotiated price lists,; account-specific discounts,
  \end{itemize}
\end{frame}

\section{Feature design for pricing models}

\begin{frame}{Feature design for pricing models}
  \begin{itemize}
    \item \textbf{Feature governance:} Pricing is high impact, so feature governance matters.
  \end{itemize}
\end{frame}

\section{Experimentation in pricing}

\begin{frame}{Experimentation in pricing}
  \begin{itemize}
    \item \textbf{Why pricing experiments are sensitive:} affect revenue directly,; may trigger competitor response,
    \item \textbf{What to measure:} conversion rate,; unit sales,
    \item \textbf{Guardrails:} margin floor protection,; inventory depletion limits,
  \end{itemize}
\end{frame}

\section{A practical pricing architecture}

\begin{frame}{A practical pricing architecture}
  \begin{itemize}
    \item \textbf{Decision flow:} gather product, demand, inventory, and market signals,; forecast demand under candidate price scenarios,
    \item \textbf{Why architecture matters:} Pricing is not just a model output.
  \end{itemize}
\end{frame}

\section{Common failure modes}

\begin{frame}{Common failure modes}
  \begin{itemize}
    \item forecasting ignores stockouts and mistakes absence of sales for lack of demand,
    \item elasticity estimates are biased by historical policy,
    \item optimization recommends prices that violate margin or contractual constraints,
    \item promotions are evaluated on revenue alone while margin collapses,
    \item competitor signals are stale or incomparable,
    \item different channels display conflicting prices due to synchronization delays.
  \end{itemize}
\end{frame}

\section{Hands-on lab: forecast plus constrained optimization}

\begin{frame}{Hands-on lab: forecast plus constrained optimization}
  \begin{itemize}
    \item \textbf{Lab scenario:} Students are given historical SKU sales, prices, promotion flags, inventory levels, and seasonality indicators.
    \item \textbf{Lab tasks:} Build a baseline demand forecast for a selected product or category.; Estimate how demand changes across a small set of candidate prices.
    \item \textbf{Expected learning outcome:} By the end of the lab, students should be able to explain why pricing requires a pipeline of assumptions, constraints, and measurement rather than a single predictive model.
  \end{itemize}
\end{frame}

\section{Managerial implications}

\begin{frame}{Managerial implications}
  \begin{itemize}
    \item clear pricing objectives,
    \item data quality around cost, stock, and catalog attributes,
    \item explicit separation between descriptive analytics and prescriptive decisions,
    \item and controlled experimentation before scaling automated pricing.
  \end{itemize}
\end{frame}

\end{document}
